\documentclass[11pt]{article}

\usepackage[a4paper,margin=1in]{geometry}
\usepackage{amsmath,amssymb}
\usepackage{algorithm}
\usepackage{algorithmic}
\usepackage{setspace}

\begin{document}

\section*{CSE400 – Fundamentals of Probability in Computing}
\textbf{Lecture 10: Randomized Min-Cut Algorithm}

\section*{1. Min-Cut Problem}

\subsection*{1.1 Why Use Min-Cut?}

The min-cut algorithm is used in various applications to solve problems related to:

\begin{itemize}
\item Network connectivity
\item Reliability
\item Optimization
\end{itemize}

\textbf{Applications Mentioned in Slides}

\textbf{Network Design}

Min cut helps in improving the efficiency of communication.

It helps in optimizing network flow.

The algorithm is used in network design to find the minimum capacity cut.

\textbf{Communication Networks}

For understanding the vulnerability of the networks to failures, minimum cut can be useful.

It helps building robust and fault-tolerant communication networks.

\textbf{VLSI Design}

In Very Large Scale Integration (VLSI) design, the algorithm is useful for partitioning circuits into smaller components.

This leads to reduced interconnectivity complexity.

Reference mentioned in slides:

Section 1.5, Application: A Randomized Min-Cut Algorithm in Probability and Computing (2nd Edition).

\section*{2. What is Min-Cut?}

\subsection*{2.1 Definition: Cut-Set}

A cut-set in a graph is:

A set of edges whose removal breaks the graph into two or more connected components.

\textbf{Step-by-Step Interpretation}

Start with a connected graph.

Remove a selected set of edges.

If the graph becomes disconnected (two or more components), that set of edges is a cut-set.

\subsection*{2.2 Definition: Minimum Cut}

Given a graph
\[
G = (V,E)
\]
with
\[
n
\]
vertices,

The minimum cut — or min-cut — problem is:

To find a minimum cardinality cut-set in
\[
G
\]

\textbf{Step-by-Step Reasoning}

Identify all possible cut-sets.

Compute the number of edges (cardinality) in each cut-set.

Select the cut-set with the smallest cardinality.

That smallest cut-set is called the minimum cut.

\subsection*{2.3 Sensitivity of Min-Cut Algorithms}

Min-cut algorithms such as Karger’s algorithm:

Are random.

Can be sensitive to the initial choice of edges.

If the algorithm happens to contract critical edges early in the process, it might find a smaller cut.

This explains why different runs of the algorithm may produce different results.

\section*{3. Edge Contraction (Main Operation)}

The main operation in the algorithm is edge contraction.

Edge contraction is:

An operation that removes an edge from a graph while simultaneously merging the two vertices.

\textbf{Contracting an Edge (u, v)}

When contracting edge
\[
(u,v)
\]

Step 1: Merge vertices $u$ and $v$ into one vertex.

Step 2: Eliminate all edges connecting $u$ and $v$.

Step 3: Retain all other edges in the graph.

\textbf{Properties After Contraction}

The new graph may have parallel edges.

The new graph has no self-loops.

(The slides visually demonstrate merging and removal of self-loops.)

\section*{4. Successful and Unsuccessful Min-Cut Runs}

\subsection*{4.1 Successful Min-Cut Run}

A successful min cut run refers to:

The success in the outcome of an algorithm designed to find the minimum cut in a graph.

From the slide illustration:

Random edges are selected.

Contractions are performed step-by-step.

The final two supernodes correspond to the correct minimum cut.

Thus, the algorithm correctly identifies the minimum cut in that execution.

\subsection*{4.2 Unsuccessful Min-Cut Run}

An unsuccessful min-cut run refers to:

An iteration of a min-cut algorithm where the algorithm fails to correctly identify the minimum cut of a given graph.

From the slide illustration:

Critical edges are contracted early.

The structure of the graph changes.

The final cut obtained is not the true minimum cut.

This demonstrates the probabilistic nature of the algorithm.

\section*{5. Max-Flow Min-Cut Theorem}

\subsection*{5.1 Statement of Theorem}

The max-flow min-cut theorem states:

“In a flow network, the maximum amount of flow passing from the source to the sink is equal to the total weight of the edges in a minimum cut.”

\subsection*{5.2 Definitions Used in Theorem}

Capacity of a cut

= The sum of the capacity of the edges in the cut that are oriented from a vertex $\in X$ to a vertex $\in Y$.

Minimum cut

= The cut in the network that has the smallest possible capacity.

Minimum cut capacity

= The capacity of the minimum cut.

Maximum flow

= The largest possible flow from source S to sink T.

(The slide includes a flow network diagram illustrating source S, sink T, capacities, and the cut partition.)

\section*{6. Deterministic Min-Cut Algorithm}

\subsection*{6.1 Stoer–Wagner Minimum Cut Algorithm}

Let $s$ and $t$ be two vertices of a graph $G$.

Let $G/\{s,t\}$ be the graph obtained by merging $s$ and $t$.

Then a minimum cut of $G$ can be obtained by taking the smaller of:

A minimum $s$–$t$ cut of $G$, and

A minimum cut of $G/\{s,t\}$.

\textbf{Justification (As Stated in Slides)}

Case 1:

There exists a minimum cut of $G$ that separates $s$ and $t$.

Then a minimum $s$–$t$ cut of $G$ is a minimum cut of $G$.

Case 2:

There is no minimum cut that separates $s$ and $t$.

Then a minimum cut of $G/\{s,t\}$ gives the required minimum cut.

Thus, one of the two cases must produce the global minimum cut.

\subsection*{6.2 Pseudocode}

\begin{algorithm}
\caption{Algorithm 1: MinimumCutPhase(G, a)}
\begin{algorithmic}
\STATE $A \leftarrow \{a\}$
\WHILE{$A \neq V$}
\STATE add to $A$ the most tightly connected vertex;
\ENDWHILE
\STATE return the cut weight as the “cut of the phase”;
\end{algorithmic}
\end{algorithm}

\begin{algorithm}
\caption{Algorithm 2: MinimumCut(G)}
\begin{algorithmic}
\WHILE{$|V| \ge 1$}
\STATE choose any $a$ from $V$;
\STATE MinimumCutPhase(G, a);
\STATE if the cut-of-the-phase is lighter than the current minimum cut then
\STATE store the cut-of-the-phase as the current minimum cut;
\STATE shrink $G$ by merging the two vertices added last;
\ENDWHILE
\STATE return the minimum cut;
\end{algorithmic}
\end{algorithm}

\subsection*{6.3 Deterministic Properties}

It always guarantees an exact minimum cut.

It may have higher time complexity for large graphs.

Time complexity of Stoer-Wagner algorithm:

\[
O(V \cdot E + V^2 \log V)
\]

where V is the number of vertices and E is the number of edges.

\section*{7. Randomized Min-Cut Algorithm}

\subsection*{7.1 Why Randomized Algorithm?}

Randomized algorithms:

Provide a probabilistic guarantee of success.

Provide a more accurate estimate of the minimum cut with fewer iterations.

Reasons listed in slides:

Efficiency

Parallelization

Approximation Guarantees

Avoidance of Worst-Case Instances

Heuristic Nature

Robustness

\section*{8. Karger’s Randomized Algorithm}

Procedure illustrated in slides:

Step 1: Pick a random edge.

Step 2: Remove it.

Step 3: Fuse its two vertices into one.

Step 4: Repeat until only two supernodes remain.

Step 5: Output the cut defined by the remaining two supernodes.

The example shown in slides demonstrates a run where random edges are picked in a way that the output cut is not minimal. The output cut differs from the actual minimum cut.

This illustrates the probabilistic behavior of the algorithm.

\section*{9. Recursive Randomized Min-Cut}

\begin{algorithm}
\caption{Algorithm 3: Recursive-Randomized-Min-Cut(G, $\alpha$)}
\begin{algorithmic}
\STATE Input:
\STATE An undirected multigraph G with n vertices, and an integer constant $\alpha > 0$.
\STATE Output:
\STATE A cut C of G.
\IF{$n \le 3$}
\STATE $C \leftarrow$ a min-cut of G found using brute force (exhaustive) search;
\ELSE
\FOR{$i \leftarrow 1$ to $\alpha$}
\STATE $G' \leftarrow$ multigraph obtained by applying $n - n/\sqrt{\alpha}$ random contraction steps on G;
\STATE $C' \leftarrow$ Recursive-Randomized-Min-Cut($G', \alpha$);
\STATE if $i = 1$ or $|C'| < |C|$ then
\STATE $C \leftarrow C'$;
\ENDFOR
\ENDIF
\STATE return C;
\end{algorithmic}
\end{algorithm}

\section*{10. Comparison: Deterministic vs Randomized Min-Cut}

Any specific problem is responsible for deciding the approach.

\textbf{Deterministic Min Cut}

Always guarantees an exact minimum cut.

May have higher time complexity for large graphs.

Time complexity:

\[
O(V \cdot E + V^2 \log V)
\]

\textbf{Randomized Min Cut}

Provides an approximate minimum cut with high probability.

Karger’s algorithm has time complexity:

\[
O(V^2)
\]

where V is the number of vertices.

\section*{11. Theorem for Min-Cut Set}

The algorithm outputs a min-cut set with probability at least:

\[
\frac{2}{n(n-1)}
\]

\end{document}
